% !TeX root = ../main.tex
\begin{chineseabstract}
% \footnotetext{*本研究得到某某基金（编号：）的资助}
% \astfootnote{本研究得到某某基金（编号：）的资助}

强化学习（Reinforcement Learning, RL）是机器学习领域的重要研究方向，通过智能体与环境的交互学习最优决策策略，已在游戏、控制等任务中取得显著成果。然而，现有算法普遍依赖理想观测条件，在真实环境部署时面临复杂观测带来的突出挑战。在视觉强化学习中，高维图像观测易受背景、光照等无关特征干扰，缺乏任务引导的数据增强在提升泛化性的同时容易破坏任务关键信息，导致样本效率低下与训练不稳定；在离线强化学习中，策略基于完全可观测数据集训练，但部署环境通常呈现部分可观测特性，状态维度因未知模式遮蔽而不可见，训练与部署间的观测差异导致价值过估计与策略性能退化。二者均严重制约强化学习在真实非平稳环境中的可靠部署。为此，本文分别针对上述两类观测挑战提出了两种算法：
\begin{enumerate}[label=\arabic*）]
    \item 针对视觉强化学习中因数据增强使用不当导致的训练不稳定与样本低效问题，提出一种任务感知型视觉数据增强选择方法（PICK）。首先，通过特征编码器提取高维观测的紧凑表征，并利用预测网络动态学习各增强方法在当前任务下的最优选用概率，构建自适应增强选择机制。随后，该机制在保留数据多样性的同时自动抑制破坏性增强操作，实现增强收益与信息保真间的平衡。实验表明，PICK在DMControlGB2基准测试的未见视觉场景中展现出优异的泛化性能，在训练稳定性与样本效率上显著优于现有先进算法。
    \item 针对离线强化学习中因部分观测缺失导致的策略泛化脆弱问题，提出一种基于自监督掩码重建的鲁棒离线强化学习方法（Masked-IQL）。首先，构建鲁棒掩码表征学习机制，引入可学习掩码标记替代零值填充以消除缺失维度的语义歧义，采用随机掩码比例采样使编码器适应广泛的缺失率范围，并通过LayerNorm与一致性损失约束残缺表征向完整表征对齐，同时融合轻量历史窗口聚合模块实现时序信息的有效利用。其次，提出两阶段解耦训练框架，将表征学习与策略优化分离，通过梯度截断消除时序差分梯度对表征结构的干扰。实验表明，Masked-IQL在多种部分可观测场景下均展现出优异的泛化性能与鲁棒性。
\end{enumerate}
\chineseKeywords{强化学习；复杂观测环境；表征优化；数据增强；离线强化学习；部分可观测决策过程}

% 关键词由3～5个组成。关键词应从《汉语主题词表》中摘选，当《汉语主题词表》的词不足以反映主题时，可由申请人设计关键词，但须加注。每一关键词之间用分号分开，最后一个关键词后不打标点符号。由申请人设计的关键词，须在该关键词的右上角标注*，并在该页的页脚处注明“*表示非汉语主题词”。

\chineseType{理论研究}

% 论文类型包括：a.理论研究(Theoretical Research)；b.应用基础(Application Fundamentals)；c.应用研究(Application Research)；d.研究报告(Research Report)；e.设计报告(Design Report)；f.案例分析(Case Study)；g.调研报告(Investigation Report)；h.产品研发(Product Development)；i.工程设计(Engineering Design)；j.工程/项目管理(Engineering/Project Management)；k.其它（Others）。

\end{chineseabstract}